\documentclass[a4paper,12pt]{article}
\usepackage[T2A]{fontenc}
\usepackage[utf8]{inputenc}
\usepackage[russian]{babel}
\usepackage{geometry}
\usepackage{xcolor}
\usepackage{tikz}
\usepackage{microtype}
\geometry{left=24mm,right=24mm,top=25mm,bottom=24mm}
\pagestyle{empty}
\definecolor{accent}{HTML}{008080}
\begin{document}
\begin{center}
\vspace*{28mm}
{\small\color{accent} МАТЕМАТИКА\par}
\vspace{15mm}
{\LARGE Ревью домашней работы\par}
\vspace{8mm}
{\Large\color{accent} Делимость и чётность\par}
\vspace{21mm}
{\large Ксения Клыкова\par}
\vspace{5mm}
16 ноября 2026 г.\par
\vfill
\begin{minipage}{0.78\linewidth}
\centering
{\large\bfseries Результат проверки\par}
\vspace{6mm}
Проверены все десять заданий. Математических ошибок,
неполных решений и существенных пропусков не обнаружено.\par
\vspace{6mm}
{\color{accent}\bfseries Рекомендация}\par
\vspace{2mm}
При доказательствах для произвольного целого числа полезно
явно выделять множитель два и остаток нуль или один.
Это помогает сделать обоснование более строгим.\par
\end{minipage}
\vfill
{\small Лёвин Артём Александрович\par}
\vspace{1mm}
{\small эксклюзивно для Ксении Клыковой\par}
\vspace{17mm}
\begin{tikzpicture}
\draw[accent,line width=0.8pt]
  (0,0) .. controls (1.1,0.55) and (2.2,-0.40) .. (3.3,0)
        .. controls (4.4,0.48) and (5.6,-0.48) .. (7.0,0)
        .. controls (8.0,0.32) and (9.0,-0.30) .. (10.0,0);
\end{tikzpicture}
\end{center}
\end{document}
